\documentclass[UTF8,12pt,a4paper,fontset=none]{ctexart}
\usepackage[a4paper,left=1.82cm,right=1.82cm,top=1.72cm,bottom=1.82cm,headheight=15pt]{geometry}
\usepackage{fontspec,xeCJK}
\setmainfont{Liberation Serif}
\setsansfont{DejaVu Sans}
\setmonofont{DejaVu Sans Mono}
\setCJKmainfont[BoldFont=Noto Serif CJK SC Bold]{Noto Serif CJK SC}
\setCJKsansfont[BoldFont=Noto Sans CJK SC Bold]{Noto Sans CJK SC}
\setCJKmonofont{Noto Sans Mono CJK SC}
\usepackage{amsmath,amssymb,mathtools,bm}
\usepackage{booktabs,longtable,tabularx,array,multirow}
\usepackage{enumitem,xcolor}
\usepackage[most]{tcolorbox}
\usepackage{fancyhdr,titlesec,hyperref,cleveref,lastpage}
\usepackage{tikz,float,caption,listings}
\usetikzlibrary{arrows.meta,positioning,calc,fit,backgrounds,shapes.geometric}
\definecolor{mainblue}{RGB}{39,82,138}
\definecolor{deepblue}{RGB}{25,55,95}
\definecolor{softblue}{RGB}{235,243,252}
\definecolor{softgreen}{RGB}{238,248,241}
\definecolor{softorange}{RGB}{255,246,232}
\definecolor{softgray}{RGB}{245,246,248}
\definecolor{warnred}{RGB}{160,48,48}
\hypersetup{unicode=true,colorlinks=true,linkcolor=deepblue,urlcolor=mainblue,citecolor=deepblue}
\setlength{\parindent}{2em}
\setlength{\parskip}{0.36em}
\setlength{\tabcolsep}{5pt}
\renewcommand{\arraystretch}{1.2}
\allowdisplaybreaks[4]
\emergencystretch=3em
\sloppy
\pagestyle{fancy}\fancyhf{}\fancyfoot[C]{\small 第 \thepage\ 页，共 \pageref{LastPage} 页}\renewcommand{\headrulewidth}{0.4pt}
\titleformat{\section}{\Large\bfseries\color{deepblue}}{\thesection}{0.65em}{}
\titleformat{\subsection}{\large\bfseries\color{mainblue}}{\thesubsection}{0.55em}{}
\titleformat{\subsubsection}{\normalsize\bfseries}{\thesubsubsection}{0.5em}{}
\numberwithin{equation}{section}
\newcommand{\R}{\mathbb{R}}
\newcommand{\E}{\mathbb{E}}
\newcommand{\Prob}{\mathbb{P}}
\newcommand{\Loss}{\mathcal{L}}
\newcommand{\norm}[1]{\left\lVert #1\right\rVert}
\newcommand{\ip}[2]{\left\langle #1,#2\right\rangle}
\newcommand{\tr}{\operatorname{Tr}}
\newcommand{\diag}{\operatorname{diag}}
\newcommand{\rank}{\operatorname{rank}}
\newcommand{\softmax}{\operatorname{softmax}}
\newcommand{\argmin}{\operatorname*{arg\,min}}
\newcommand{\argmax}{\operatorname*{arg\,max}}
\newcommand{\sg}{\operatorname{sg}}
\newcommand{\KL}{D_{\mathrm{KL}}}
\newcommand{\dd}{\mathrm{d}}
\newtcolorbox{keybox}[1][]{enhanced,breakable,colback=softblue,colframe=mainblue,boxrule=0.8pt,arc=1.5mm,left=2mm,right=2mm,top=1.2mm,bottom=1.2mm,title={#1},fonttitle=\bfseries}
\newtcolorbox{examplebox}[1][]{enhanced,breakable,colback=softgreen,colframe=green!45!black,boxrule=0.7pt,arc=1.5mm,left=2mm,right=2mm,top=1.2mm,bottom=1.2mm,title={#1},fonttitle=\bfseries}
\newtcolorbox{notebox}[1][]{enhanced,breakable,colback=softorange,colframe=orange!70!black,boxrule=0.7pt,arc=1.5mm,left=2mm,right=2mm,top=1.2mm,bottom=1.2mm,title={#1},fonttitle=\bfseries}
\newtcolorbox{criticalbox}[1][]{enhanced,breakable,colback=red!3,colframe=warnred,boxrule=0.8pt,arc=1.5mm,left=2mm,right=2mm,top=1.2mm,bottom=1.2mm,title={#1},fonttitle=\bfseries}
\lstset{basicstyle=\ttfamily\small,breaklines=true,frame=single,backgroundcolor=\color{softgray},numbers=left,numberstyle=\tiny,showstringspaces=false,columns=fullflexible}
\hypersetup{pdftitle={04 NoProp 数学过程详解},pdfauthor={OpenAI}}
\fancyhead[L]{\small 04 NoProp}
\fancyhead[R]{\small 数学过程详解}
\setlength{\abovedisplayskip}{0.82em}
\setlength{\belowdisplayskip}{0.82em}
\setlength{\abovedisplayshortskip}{0.45em}
\setlength{\belowdisplayshortskip}{0.45em}
\begin{document}
\begin{center}
{\LARGE\bfseries 04\_NoProp 数学过程详解}\par
\vspace{0.35em}
{\large ELBO、高斯条件分布、SNR 权重与局部独立训练的完整推导}
\end{center}
\vspace{0.45em}
\begin{keybox}[本文只处理真正需要推导的部分]
本文重点还原论文附录 A.1--A.4：从证据下界到高斯 KL，再到 SNR 差分加权的逐层平方误差，并解释连续时间和 Flow Matching 版本。全文按“公式从哪里来--每个符号是什么--中间步骤为何成立--维度和复杂度如何核对--论文结论依赖哪些假设”的顺序展开；不复述实验表格，也不使用与论文无关的通用背景凑页数。
\end{keybox}
\tableofcontents
\newpage

\section{NoProp 的概率模型：把网络深度改写成潜变量链}
论文令类别标签 $y\in\{1,\ldots,m\}$ 对应嵌入 $u_y\in\R^d$。深度方向的隐状态记为 $z_0,z_1,\ldots,z_T$。生成方向（论文称 forward process）为
\begin{equation}
 p(z_{0:T},y\mid x)=p(z_0)\prod_{t=1}^{T}p_\theta(z_t\mid z_{t-1},x)\,p_{\theta_{\rm out}}(y\mid z_T).
\end{equation}
辅助后验方向为
\begin{equation}
 q(z_{0:T}\mid y,x)=q(z_T\mid y)\prod_{t=T}^{1}q(z_{t-1}\mid z_t,y).
\end{equation}
这里 $q$ 只依赖标签嵌入并被固定成可解析高斯扩散；$p_\theta$ 由每一层独立网络块拟合。NoProp 的“无全局前向/反向传播”来自：每个 $t$ 的训练样本 $z_{t-1}$ 可以直接从 $q(z_{t-1}\mid y)$ 采样，而不必先运行前面所有块。

\section{ELBO 从哪里来：逐步恢复论文公式}
从边缘似然
\begin{equation}
 p(y\mid x)=\int p(z_{0:T},y\mid x)\,\dd z_{0:T}
\end{equation}
出发，乘除任意后验 $q(z_{0:T}\mid y,x)$：
\begin{align}
 \log p(y\mid x)
 &=\log\int q(z_{0:T}\mid y,x)
 \frac{p(z_{0:T},y\mid x)}{q(z_{0:T}\mid y,x)}\,\dd z_{0:T}.
\end{align}
由于 $\log$ 为凹函数，Jensen 不等式给出
\begin{equation}
 \boxed{\log p(y\mid x)\ge
 \E_q\left[\log p(z_{0:T},y\mid x)-\log q(z_{0:T}\mid y,x)\right].}
\end{equation}
把两个链式分解代入并重排，可得
\begin{align}
 \operatorname{ELBO}
 =&\;\E_{q(z_T\mid y)}\log p_{\theta_{\rm out}}(y\mid z_T)
 -\KL\!\left(q(z_0\mid y)\Vert p(z_0)\right)\\
 &-\sum_{t=1}^{T}\E_{q(z_{t-1}\mid y)}
 \KL\!\left(q(z_t\mid z_{t-1},y)\Vert p_\theta(z_t\mid z_{t-1},x)\right).
\end{align}
因此最小化负 ELBO 会产生三个部分：输出交叉熵、初始分布 KL、每层条件高斯 KL。

\section{固定 OU 扩散后验的边缘分布}
论文使用方差保持的离散 Ornstein--Uhlenbeck 过程。设单步系数 $\alpha_t\in(0,1)$，累乘
\begin{equation}
 \bar\alpha_t=\prod_{s=t+1}^{T}\alpha_s
\end{equation}
（索引方向按论文由 $z_T$ 向 $z_0$ 的后验链约定）。边缘分布写成
\begin{equation}
 \boxed{q(z_t\mid y)=\mathcal N\!\left(z_t;\sqrt{\bar\alpha_t}u_y,(1-\bar\alpha_t)I\right).}
\end{equation}
所以可直接重参数化采样
\begin{equation}
 z_t=\sqrt{\bar\alpha_t}u_y+\sqrt{1-\bar\alpha_t}\,\epsilon,
 \qquad \epsilon\sim\mathcal N(0,I).
\end{equation}
均值保留标签信号，方差注入噪声。信噪比为
\begin{equation}
 \operatorname{SNR}(t)=\frac{\bar\alpha_t}{1-\bar\alpha_t}.
\end{equation}

\section{条件高斯均值 $a_tu_y+b_tz_{t-1}$ 的推导思想}
联合变量 $(z_{t-1},z_t)$ 在给定 $y$ 时为高斯，因此条件分布仍为高斯：
\begin{equation}
 q(z_t\mid z_{t-1},y)=\mathcal N\!\left(z_t;\mu_t(z_{t-1},u_y),c_tI\right).
\end{equation}
高斯条件均值公式
\begin{equation}
 \E[z_t\mid z_{t-1},y]
 =\mu_t^{(0)}+\Sigma_{t,t-1}\Sigma_{t-1,t-1}^{-1}(z_{t-1}-\mu_{t-1}^{(0)})
\end{equation}
说明条件均值必为 $z_{t-1}$ 与 $u_y$ 的线性组合，因此可写成
\begin{equation}
 \mu_t(z_{t-1},u_y)=a_tu_y+b_tz_{t-1}.
\end{equation}
论文给出的系数为
\begin{align}
 a_t&=\frac{\sqrt{\bar\alpha_t}(1-\alpha_{t-1})}{1-\bar\alpha_{t-1}},\\
 b_t&=\frac{\sqrt{\alpha_{t-1}}(1-\bar\alpha_t)}{1-\bar\alpha_{t-1}},\\
 c_t&=\frac{(1-\bar\alpha_t)(1-\alpha_{t-1})}{1-\bar\alpha_{t-1}}.
\end{align}
这些系数并非可学习参数，而是由噪声日程唯一决定。

\section{网络转移与残差结构}
模型把真实标签嵌入 $u_y$ 替换成网络预测 $\hat u_{\theta_t}(z_{t-1},x)$：
\begin{equation}
 p_\theta(z_t\mid z_{t-1},x)
 =\mathcal N\!\left(z_t;
 a_t\hat u_{\theta_t}(z_{t-1},x)+b_tz_{t-1},c_tI\right).
\end{equation}
重参数化后
\begin{equation}
 \boxed{z_t=a_t\hat u_{\theta_t}(z_{t-1},x)+b_tz_{t-1}+\sqrt{c_t}\epsilon_t.}
\end{equation}
$b_tz_{t-1}$ 是显式跳连，$a_t\hat u_{\theta_t}$ 是当前块对标签原型的估计，噪声项保证与扩散后验方差一致。

\section{高斯 KL 如何化为逐层平方误差}
两个协方差完全相同的高斯
\begin{equation}
 q=\mathcal N(\mu_q,c_tI),\qquad p=\mathcal N(\mu_p,c_tI)
\end{equation}
满足
\begin{equation}
 \KL(q\Vert p)=\frac{1}{2c_t}\norm{\mu_q-\mu_p}_2^2.
\end{equation}
本论文中
\begin{align}
 \mu_q&=a_tu_y+b_tz_{t-1},\\
 \mu_p&=a_t\hat u_{\theta_t}(z_{t-1},x)+b_tz_{t-1}.
\end{align}
跳连项相消：
\begin{equation}
 \mu_q-\mu_p=a_t\left(u_y-\hat u_{\theta_t}\right).
\end{equation}
于是
\begin{equation}
 \KL(q\Vert p)=\frac{a_t^2}{2c_t}
orm{\hat u_{\theta_t}-u_y}_2^2.
\end{equation}
代入 $a_t,c_t$ 并化简：
\begin{align}
 \frac{a_t^2}{2c_t}
 &=\frac12\frac{\bar\alpha_t(1-\alpha_{t-1})}
 {(1-\bar\alpha_t)(1-\bar\alpha_{t-1})}\\
 &=\frac12\left[
 \frac{\bar\alpha_t}{1-\bar\alpha_t}
 -\frac{\bar\alpha_{t-1}}{1-\bar\alpha_{t-1}}
 \right]\\
 &=\frac12\left[\operatorname{SNR}(t)-\operatorname{SNR}(t-1)\right].
\end{align}
所以每层 KL 正好变成 SNR 增量加权的原型回归损失。

\section{NoProp 离散目标的完整形式}
取负 ELBO，并用均匀采样 $t\sim U\{1,\ldots,T\}$ 对求和做无偏估计，得到
\begin{align}
 \Loss_{\mathrm{NoProp}}
 =&\;\E_{q(z_T\mid y)}[-\log \hat p_{\theta_{\rm out}}(y\mid z_T)]
 +\KL(q(z_0\mid y)\Vert p(z_0))\\
 &+\frac{\eta T}{2}\E_{t,z_{t-1}}
 \left[(\operatorname{SNR}(t)-\operatorname{SNR}(t-1))
 \norm{\hat u_{\theta_t}(z_{t-1},x)-u_y}_2^2\right].
\end{align}
若实现中遍历所有 $t$，则不需要 $T$ 的 Monte Carlo 补偿；若随机抽一个 $t$ 估计总和，则必须乘 $T$ 才严格无偏。论文公式将常数吸收到超参数 $\eta$ 中。

\section{为什么每个块可以独立训练}
第 $t$ 块的局部损失为
\begin{equation}
 \ell_t(\theta_t)=w_t\norm{\hat u_{\theta_t}(z_{t-1},x)-u_y}_2^2,
 \qquad w_t=\operatorname{SNR}(t)-\operatorname{SNR}(t-1).
\end{equation}
训练输入 $z_{t-1}$ 直接从已知 $q(z_{t-1}\mid y)$ 采样，所以
\begin{equation}
 \nabla_{\theta_t}\ell_t
 =2w_tJ_{\theta_t}\hat u_{\theta_t}^{\top}(\hat u_{\theta_t}-u_y)
\end{equation}
不包含 $\theta_1,\ldots,\theta_{t-1}$ 的导数。标准反向传播则需要
\begin{equation}
 \frac{\partial\Loss}{\partial\theta_t}
 =\frac{\partial\Loss}{\partial z_T}
 \prod_{s=t+1}^{T}\frac{\partial z_s}{\partial z_{s-1}}
 \frac{\partial z_t}{\partial\theta_t}.
\end{equation}
NoProp 用固定后验提供局部监督，避免了这个长 Jacobian 乘积，但代价是优化目标变成变分下界而非直接端到端分类损失。

\section{连续时间极限}
令时间间隔 $\tau=1/T$。离散权重可写成差商
\begin{equation}
 \frac{\operatorname{SNR}(t+\tau)-\operatorname{SNR}(t)}{\tau}.
\end{equation}
当 $T\to\infty$、$\tau\to0$ 时
\begin{equation}
 \frac{\operatorname{SNR}(t+\tau)-\operatorname{SNR}(t)}{\tau}
 \to \operatorname{SNR}'(t).
\end{equation}
连续目标为
\begin{align}
 \Loss_{\mathrm{CT}}
 =&\;\E[-\log\hat p(y\mid z_1)]
 +\KL(q(z_0\mid y)\Vert p(z_0))\\
 &+\frac{\eta}{2}\E_{t\sim U[0,1]}
 \left[\operatorname{SNR}'(t)
 \norm{\hat u_\theta(z_t,x,t)-u_y}_2^2\right].
\end{align}
离散版本每个时刻一个块；连续版本用同一网络并把 $t$ 作为输入。

\section{Flow Matching 版本}
取 $z_0\sim\mathcal N(0,I)$、$z_1=u_y$，定义高斯路径
\begin{equation}
 p_t(z_t\mid z_0,z_1,x)=\mathcal N\!\left(z_t;t z_1+(1-t)z_0,\sigma^2I\right).
\end{equation}
均值轨迹的导数为
\begin{equation}
 \frac{\dd}{\dd t}\left[t z_1+(1-t)z_0\right]=z_1-z_0.
\end{equation}
故条件向量场是常量 $z_1-z_0$，训练目标
\begin{equation}
 \boxed{\Loss_{\mathrm{FM}}=
 \E\norm{v_\theta(z_t,x,t)-(u_y-z_0)}_2^2.}
\end{equation}
与离散 NoProp 相同，它也允许随机抽取时间训练，不需要把整条 ODE 路径展开后反向传播。

\section{输出分类器与标签嵌入重构}
最后输出 logits $f_{\theta_{\rm out}}(z_T)\in\R^m$，概率
\begin{equation}
 \hat p(y\mid z_T)=\frac{\exp f_y(z_T)}{\sum_{j=1}^{m}\exp f_j(z_T)}.
\end{equation}
若将概率重新映射回嵌入空间，
\begin{equation}
 \tilde y=\softmax(f(z_T))W_{\rm Embed},
\end{equation}
则 $\tilde y$ 是所有类别原型的凸组合。平方重构
\begin{equation}
 \norm{\tilde y-u_y}_2^2
\end{equation}
不仅惩罚分类错误，还利用类别嵌入之间的几何距离；若原型不正交，两个错误类别的惩罚可以不同。

\section{NoProp 结论的严格边界}
\begin{criticalbox}[“不需要前向传播”不等于推理不前向]
训练第 $t$ 块时不需要先运行 $1,\ldots,t-1$ 块，也不需要把最终损失反传穿过全部深度；但推理时仍需按 $z_0\to z_1\to\cdots\to z_T$ 运行网络链。局部目标来自固定后验与 ELBO，因此其最优解是否达到端到端训练性能取决于后验设计、标签嵌入容量和各局部块的拟合误差。
\end{criticalbox}
% AUGMENTED_DETAILED_MATH

\section{ELBO 的逐步 Jensen 推导}
观测标签为 $y$，潜变量链为 $z_{0:T}$。边缘似然
\begin{equation}
 p_\theta(y\mid x)=\int p_\theta(y,z_{0:T}\mid x)\dd z_{0:T}.
\end{equation}
插入任意变分分布 $q(z_{0:T}\mid y)$：
\begin{align}
 \log p_\theta(y\mid x)
 &=\log\E_q\left[
 \frac{p_\theta(y,z_{0:T}\mid x)}{q(z_{0:T}\mid y)}
 \right]\\
 &\ge\E_q\log
 \frac{p_\theta(y,z_{0:T}\mid x)}{q(z_{0:T}\mid y)}
 =\mathcal E(\theta),
\end{align}
不等号来自 $\log$ 的凹性。两者差值恰好是
\begin{equation}
 \log p_\theta(y\mid x)-\mathcal E
 =\KL(q(z_{0:T}\mid y)\Vert p_\theta(z_{0:T}\mid y,x))\ge0.
\end{equation}
因此最大化 ELBO 等价于让模型后验逼近固定的标签扩散后验。

\section{Markov 分解如何产生逐层 KL}
若
\begin{equation}
 q(z_{0:T}\mid y)=q(z_0\mid y)\prod_{t=1}^{T}q(z_t\mid z_{t-1},y)
\end{equation}
且生成模型
\begin{equation}
 p_\theta(y,z_{0:T}\mid x)
 =p(y\mid z_T)p(z_0)\prod_{t=1}^{T}p_{\theta_t}(z_t\mid z_{t-1},x),
\end{equation}
则对数比值拆成各因子的对数之和。取期望后
\begin{align}
 \mathcal E
 &=\E_q\log p(y\mid z_T)
 -\KL(q(z_0\mid y)\Vert p(z_0))\\
 &\quad-\sum_{t=1}^{T}
 \E_{q(z_{t-1}\mid y)}
 \KL(q(z_t\mid z_{t-1},y)
 \Vert p_{\theta_t}(z_t\mid z_{t-1},x)).
\end{align}
每个 $\theta_t$ 只出现在第 $t$ 个 KL 中，这是可局部训练的概率根源。

\section{OU 边缘分布的 SDE 解}
连续 OU 过程
\begin{equation}
 \dd z_t=-\frac12\beta(t)z_t\dd t+\sqrt{\beta(t)}\dd W_t,
 \qquad z_0=u_y
\end{equation}
的积分因子为
\begin{equation}
 a(t)=\exp\left(-\frac12\int_0^t\beta(s)\dd s\right).
\end{equation}
解为
\begin{equation}
 z_t=a(t)u_y+
 a(t)\int_0^t\frac{\sqrt{\beta(s)}}{a(s)}\dd W_s.
\end{equation}
随机积分是零均值高斯，方差
\begin{align}
 a(t)^2\int_0^t\frac{\beta(s)}{a(s)^2}\dd s
 &=a(t)^2[a(t)^{-2}-1]\\
 &=1-a(t)^2.
\end{align}
故
\begin{equation}
 \boxed{q(z_t\mid y)=\mathcal N(a(t)u_y,(1-a(t)^2)I).}
\end{equation}
离散记号中 $a(t)^2$ 对应 $\bar\alpha_t$。

\section{高斯条件后验的配方法}
设
\begin{equation}
 q(z_{t-1}\mid y)=\mathcal N(m_{t-1},s_{t-1}^2I),
\end{equation}
转移
\begin{equation}
 q(z_t\mid z_{t-1})=\mathcal N(\sqrt{\alpha_t}z_{t-1},(1-\alpha_t)I).
\end{equation}
Bayes 后验的负对数中与 $z_{t-1}$ 有关的项为
\begin{equation}
 \frac{\norm{z_{t-1}-m_{t-1}}^2}{2s_{t-1}^2}
 +\frac{\norm{z_t-\sqrt{\alpha_t}z_{t-1}}^2}{2(1-\alpha_t)}.
\end{equation}
展开并收集二次项，精度矩阵
\begin{equation}
 \Lambda=\frac1{s_{t-1}^2}I+\frac{\alpha_t}{1-\alpha_t}I.
\end{equation}
均值为
\begin{equation}
 \mu=\Lambda^{-1}
 \left(\frac{m_{t-1}}{s_{t-1}^2}
 +\frac{\sqrt{\alpha_t}}{1-\alpha_t}z_t\right).
\end{equation}
代入 $m_{t-1}=\sqrt{\bar\alpha_{t-1}}u_y$ 与 $s_{t-1}^2=1-\bar\alpha_{t-1}$，即可整理成论文的
\begin{equation}
 \mu_t=a_tu_y+b_tz_t,
\end{equation}
并得到闭式方差 $c_tI$。

\section{等协方差高斯 KL 的完整化简}
一般高斯 KL 为
\begin{align}
 \KL(\mathcal N(\mu_q,\Sigma_q)\Vert
 \mathcal N(\mu_p,\Sigma_p))
 &=\frac12\left[
 \tr(\Sigma_p^{-1}\Sigma_q)
 +(\mu_p-\mu_q)^\top\Sigma_p^{-1}(\mu_p-\mu_q)
 \right.\\
 &\quad\left.-d+\log\frac{\det\Sigma_p}{\det\Sigma_q}
 \right].
\end{align}
当 $\Sigma_q=\Sigma_p=c_tI$ 时，迹项为 $d$，对数行列式为 0，与 $-d$ 抵消，只剩
\begin{equation}
 \KL=\frac1{2c_t}\norm{\mu_q-\mu_p}^2.
\end{equation}
若两均值只在标签嵌入预测上不同
\begin{equation}
 \mu_q-\mu_p=a_t(u_y-\widehat u_{\theta_t}),
\end{equation}
则
\begin{equation}
 \boxed{
 \KL=\frac{a_t^2}{2c_t}
orm{u_y-\widehat u_{\theta_t}}^2.
 }
\end{equation}
这说明论文逐层 MSE 不是随意设计，而是精确高斯 KL。

\section{局部目标为何不需要跨层反向传播}
第 $t$ 层损失
\begin{equation}
 \Loss_t(\theta_t)=
 \E_{y,x,z_{t-1}}
 w_t\norm{u_y-f_{\theta_t}(z_{t-1},x)}^2.
\end{equation}
采样 $z_{t-1}$ 来自固定 $q$，不依赖任何前层参数。因此
\begin{equation}
 \frac{\partial z_{t-1}}{\partial\theta_s}=0,
 \qquad s<t.
\end{equation}
于是
\begin{equation}
 \nabla_{\theta_s}\Loss_t=0\quad(s\ne t),
\end{equation}
每个块可独立优化。与 backprop 的差别不是链式法则被“近似”，而是概率构造主动切断了参数依赖图。

\section{连续时间 Flow Matching 目标}
在连续时间中，标签条件路径 $q_t(z\mid y)$ 满足连续性方程
\begin{equation}
 \partial_tq_t+\nabla\cdot(q_tu_t^y)=0.
\end{equation}
模型场 $v_\theta(z,t,x)$ 用条件平方损失学习
\begin{equation}
 \Loss_{\mathrm{FM}}
 =\E_{t,y,z_t}\norm{v_\theta(z_t,t,x)-u_t^y(z_t)}^2.
\end{equation}
对固定 $(z,t,x)$，最优模型同样是条件期望
\begin{equation}
 v^*(z,t,x)=\E[u_t^y(z)\mid z_t=z,x].
\end{equation}
因此离散高斯 KL 版与连续 FM 版都在学习“由输入 $x$ 预测标签扩散方向”，只是一者通过转移密度，一者通过速度场表达。

\section{内存复杂度对比}
普通 $T$ 层网络反向传播需保存各层激活，内存近似
\begin{equation}
 M_{\mathrm{BP}}=O(TBd),
\end{equation}
其中 $B$ 为 batch，$d$ 为激活宽度。NoProp 单独训练某块时只需该块输入与局部中间量
\begin{equation}
 M_{\mathrm{NP}}=O(Bd)+M_{\mathrm{block}}.
\end{equation}
但总计算并不自动减少：若每个块都需要单独采样和更新，总训练 FLOPs 仍与 $T$ 近似线性。其主要收益是并行性、局部优化和激活内存，而不是“免费训练”。

\section{输出分类器与标签嵌入几何}
若最终表示 $z_T$ 与类别嵌入 $u_c$ 用距离分类
\begin{equation}
 p(y=c\mid z_T)
 =\frac{\exp(-\norm{z_T-u_c}^2/\tau)}
 {\sum_j\exp(-\norm{z_T-u_j}^2/\tau)},
\end{equation}
展开平方距离
\begin{equation}
 -\norm{z-u_c}^2
 =2z^\top u_c-\norm{u_c}^2-\norm z^2.
\end{equation}
对所有类别公共的 $-\norm z^2$ 在 softmax 中抵消，故等价于带类别偏置的线性分类器
\begin{equation}
 \operatorname{logit}_c=\frac{2}{\tau}z^\top u_c-\frac1\tau\norm{u_c}^2.
\end{equation}
\clearpage
\section{从联合分布到 ELBO 的逐行推导}
设潜变量链为 $z_{0:T}$，观测标签为 $y$，输入 $x$ 作为条件。模型联合分布
\begin{equation}
 p_\theta(y,z_{0:T}\mid x)
 =p(z_0)\prod_{t=1}^{T}p_\theta(z_t\mid z_{t-1},x)\,p_\theta(y\mid z_T).
\end{equation}
引入任意后验近似 $q(z_{0:T}\mid x,y)$，有
\begin{align}
 \log p_\theta(y\mid x)
 &=\log\int q(z_{0:T}\mid x,y)
 \frac{p_\theta(y,z_{0:T}\mid x)}{q(z_{0:T}\mid x,y)}\,\dd z_{0:T}\\
 &=\log\E_q\left[
 \frac{p_\theta(y,z_{0:T}\mid x)}{q(z_{0:T}\mid x,y)}\right].
\end{align}
由 $\log$ 的凹性与 Jensen 不等式，
\begin{equation}
 \log p_\theta(y\mid x)
 \ge \E_q\left[
 \log p_\theta(y,z_{0:T}\mid x)-\log q(z_{0:T}\mid x,y)
 \right].
\end{equation}
右侧即 ELBO。差距恰好为
\begin{equation}
 \log p_\theta(y\mid x)-\mathrm{ELBO}
 =\KL\bigl(q(z_{0:T}\mid x,y)\|p_\theta(z_{0:T}\mid x,y)\bigr)\ge0.
\end{equation}
所以最大化 ELBO 等价于同时提高数据似然并使近似后验接近真实后验。

\section{反向 Markov 分解如何产生逐层 KL}
论文选择反向扩散后验
\begin{equation}
 q(z_{0:T}\mid y)=q(z_T\mid y)\prod_{t=1}^{T}q(z_{t-1}\mid z_t,y).
\end{equation}
把它与前向模型相除并取对数：
\begin{align}
 \log\frac{q(z_{0:T}\mid y)}{p_\theta(y,z_{0:T}\mid x)}
 &=\log\frac{q(z_T\mid y)}{p_\theta(y\mid z_T)}
 +\sum_{t=1}^{T}
 \log\frac{q(z_{t-1}\mid z_t,y)}{p_\theta(z_t\mid z_{t-1},x)}
 -\log p(z_0).
\end{align}
为了得到标准扩散式逐步 KL，需要利用 Bayes 关系把 $q(z_{t-1}\mid z_t,y)$ 与前向条件写到同一方向。整理后每一层都出现一个只涉及该层转移参数的 KL 或平方误差项。这是 NoProp 能局部训练的概率根源，而不是简单地“截断反向传播”。

\section{线性高斯 OU 边缘分布的递推证明}
考虑
\begin{equation}
 z_t=\alpha_t z_{t-1}+\sqrt{1-\alpha_t^2}\,\epsilon_t,
 \qquad \epsilon_t\sim\mathcal N(0,I).
\end{equation}
假设
\begin{equation}
 z_{t-1}\mid y\sim
 \mathcal N(\bar\alpha_{t-1}u_y,(1-\bar\alpha_{t-1}^2)I),
\end{equation}
其中 $\bar\alpha_t=\prod_{s=1}^{t}\alpha_s$。线性高斯变换给出均值
\begin{equation}
 \E[z_t\mid y]=\alpha_t\bar\alpha_{t-1}u_y=\bar\alpha_tu_y,
\end{equation}
方差
\begin{align}
 \operatorname{Var}(z_t\mid y)
 &=\alpha_t^2(1-\bar\alpha_{t-1}^2)I+(1-\alpha_t^2)I\\
 &=(1-\alpha_t^2\bar\alpha_{t-1}^2)I
 =(1-\bar\alpha_t^2)I.
\end{align}
由数学归纳法，
\begin{equation}
 \boxed{z_t\mid y\sim
 \mathcal N(\bar\alpha_tu_y,(1-\bar\alpha_t^2)I).}
\end{equation}

\section{两个高斯相乘得到条件后验的配方法}
标量情形先看清结构。若
\begin{equation}
 q(z_t\mid z_{t-1})=\mathcal N(az_{t-1},\sigma^2),
 \qquad
 q(z_{t-1}\mid y)=\mathcal N(m,s^2),
\end{equation}
则
\begin{align}
 \log q(z_{t-1}\mid z_t,y)
 &=-\frac{(z_t-az_{t-1})^2}{2\sigma^2}
 -\frac{(z_{t-1}-m)^2}{2s^2}+\mathrm{const}.
\end{align}
收集 $z_{t-1}^2$ 和一次项：
\begin{equation}
 -\frac12\left(\frac{a^2}{\sigma^2}+\frac1{s^2}\right)z_{t-1}^2
 +\left(\frac{az_t}{\sigma^2}+\frac{m}{s^2}\right)z_{t-1}.
\end{equation}
因此后验方差与均值为
\begin{equation}
 \tilde s^2=\left(\frac{a^2}{\sigma^2}+\frac1{s^2}\right)^{-1},
 \qquad
 \tilde m=\tilde s^2\left(\frac{az_t}{\sigma^2}+\frac{m}{s^2}\right).
\end{equation}
多维各向同性情形逐维相同。论文中的系数 $a_t,b_t,c_t$ 正是把该后验均值写成
\begin{equation}
 \tilde m_t=a_tu_y+b_tz_t
\end{equation}
后的闭式系数。

\section{等协方差高斯 KL 为什么变成平方误差}
对
\begin{equation}
 q=\mathcal N(\mu_q,\Sigma),\qquad
 p=\mathcal N(\mu_p,\Sigma),
\end{equation}
一般高斯 KL 为
\begin{align}
 \KL(q\|p)=\frac12\Bigl[
 \tr(\Sigma^{-1}\Sigma)
 +(\mu_p-\mu_q)^\top\Sigma^{-1}(\mu_p-\mu_q)
 -d+\log\frac{\det\Sigma}{\det\Sigma}
 \Bigr].
\end{align}
迹项等于 $d$，行列式项为零，于是
\begin{equation}
 \KL(q\|p)=\frac12
 \|\mu_p-\mu_q\|_{\Sigma^{-1}}^2.
\end{equation}
若 $\Sigma=c_tI$，则
\begin{equation}
 \KL(q\|p)=\frac1{2c_t}\|\mu_p-\mu_q\|_2^2.
\end{equation}
当两者均值只在 $u_y$ 与网络预测 $\hat u_{\theta_t}$ 上不同，便得到加权局部回归损失
\begin{equation}
 \Loss_t=\lambda_t\|\hat u_{\theta_t}(z_{t-1},x)-u_y\|_2^2.
\end{equation}

\section{局部梯度为何不包含其他块参数}
设第 $t$ 块损失为
\begin{equation}
 \Loss_t(\theta_t)=\E_{z_{t-1}\sim q}
 \|u_{\theta_t}(z_{t-1},x)-u_y\|^2.
\end{equation}
关键是 $z_{t-1}$ 来自固定的教师后验 $q$，不是由前面可训练块生成。于是对 $s\ne t$，
\begin{equation}
 \frac{\partial\Loss_t}{\partial\theta_s}=0.
\end{equation}
对本块
\begin{equation}
 \nabla_{\theta_t}\Loss_t
 =2\E J_{\theta_t}u_{\theta_t}^{\top}
 (u_{\theta_t}-u_y).
\end{equation}
如果改成把模型自己的 $z_{t-1}(\theta_{<t})$ 喂给第 $t$ 块，则链式法则产生
\begin{equation}
 \frac{\partial\Loss_t}{\partial\theta_s}
 =\frac{\partial\Loss_t}{\partial z_{t-1}}
 \frac{\partial z_{t-1}}{\partial\theta_s},
\end{equation}
局部独立性立即消失。因此 NoProp 的核心是固定噪声教师分布，而非仅仅调用 stop-gradient。

\section{一个 $T=2$ 的完整标量例子}
令标签嵌入 $u_y=2$，$\alpha_1=\alpha_2=0.8$。则
\begin{equation}
 \bar\alpha_1=0.8,
 \qquad
 \bar\alpha_2=0.64.
\end{equation}
边缘分布为
\begin{equation}
 z_1\mid y\sim\mathcal N(1.6,1-0.8^2),
\end{equation}
\begin{equation}
 z_2\mid y\sim\mathcal N(1.28,1-0.64^2).
\end{equation}
训练第一块时直接采 $z_0$ 或对应教师变量，回归目标始终是 $2$；训练第二块时独立采 $z_1$，仍回归 $2$。两个优化问题
\begin{equation}
 \min_{\theta_1}\E(u_{\theta_1}(z_0,x)-2)^2,
 \qquad
 \min_{\theta_2}\E(u_{\theta_2}(z_1,x)-2)^2
\end{equation}
可以并行求梯度。推理时才把第一块输出生成的状态传给第二块，因此训练分布与推理分布可能存在偏移，这是局部训练换来的主要代价。

\section{连续时间极限：离散 OU 到 SDE}
令步长 $\Delta t$ 很小，取
\begin{equation}
 \alpha_t=1-\frac12\beta(t)\Delta t+o(\Delta t).
\end{equation}
噪声尺度
\begin{align}
 \sqrt{1-\alpha_t^2}
 &=\sqrt{1-\left(1-\beta(t)\Delta t+o(\Delta t)\right)}\\
 &=\sqrt{\beta(t)\Delta t}+o(\sqrt{\Delta t}).
\end{align}
离散更新
\begin{equation}
 z_{t+\Delta t}-z_t
 =-\frac12\beta(t)z_t\Delta t
 +\sqrt{\beta(t)}\sqrt{\Delta t}\,\epsilon
\end{equation}
在极限下成为方差保持 SDE
\begin{equation}
 \dd z_t=-\frac12\beta(t)z_t\dd t+\sqrt{\beta(t)}\dd W_t.
\end{equation}
其均值衰减因子为
\begin{equation}
 \bar\alpha(t)=\exp\left[-\frac12\int_0^t\beta(s)\dd s\right].
\end{equation}

\section{Flow Matching 版本的条件路径与速度}
若取桥接路径
\begin{equation}
 z_t=(1-t)z_0+t u_y+\sigma\epsilon,
\end{equation}
对 $t$ 求导得到条件速度
\begin{equation}
 v^*(z_t\mid z_0,u_y)=u_y-z_0.
\end{equation}
平方损失的最优向量场为条件期望
\begin{equation}
 v_\theta^*(z,t,x)=\E[u_y-z_0\mid z_t=z,x].
\end{equation}
由估计终点
\begin{equation}
 \widetilde u_y=z_t+(1-t)v_\theta(z_t,t,x),
\end{equation}
可以直接接分类器。若 $v_\theta$ 精确，代入路径有
\begin{equation}
 z_t+(1-t)(u_y-z_0)=t u_y+(1-t)z_0+(1-t)u_y-(1-t)z_0=u_y
\end{equation}
（忽略固定桥噪声时），说明该终点估计公式的系数来自直线运动几何。

\section{内存复杂度的严格比较}
标准端到端反向传播需要保存所有层激活。若每层激活大小为 $A_t$，峰值激活内存近似
\begin{equation}
 M_{\rm BP}=\Theta\left(\sum_{t=1}^{T}A_t\right).
\end{equation}
NoProp 每次只训练一个块或并行独立块，不需要保留跨块计算图，单设备顺序训练时
\begin{equation}
 M_{\rm NP}=\Theta\left(\max_t A_t\right).
\end{equation}
但总前向计算仍包含所有块，且为每层构造教师噪声样本和局部优化器状态。若所有块同时放在多设备上并行训练，总系统内存未必更小，只是单设备的反向图深度被消除。
\end{document}
